\documentclass[11pt]{article}
\usepackage{amsmath,amssymb,mathtools}
\usepackage[a4paper,margin=1in]{geometry}

\begin{document}

\section*{“Phase I” method for initial feasible point (0.25 pt)}

We consider the original constraints:
\[
\begin{aligned}
\text{(C1)}\;& x - 2y + 2 \ge 0,\\
\text{(C2)}\;& -x - 2y + 6 \ge 0,\\
\text{(C3)}\;& -x + 2y + 2 \ge 0,\\
\text{(C4)}\;& x \ge 0,\\
\text{(C5)}\;& y \ge 0.
\end{aligned}
\]

Given the starting point
\[
x_{0} = \begin{pmatrix}1\\[0.2em]2.5\end{pmatrix},
\]
we first check feasibility:
\[
\begin{aligned}
\text{C1: } &1 - 2\cdot 2.5 + 2 = -2 \quad (\text{violated by }2),\\
\text{C2: } &-1 - 2\cdot 2.5 + 6 = 0 \quad (\text{active}),\\
\text{C3: } &-1 + 2\cdot 2.5 + 2 = 6 \quad (>0),\\
\text{C4: } &1 \ge 0,\qquad
\text{C5: } 2.5 \ge 0.
\end{aligned}
\]
Only Constraint (C1) is violated.

\subsection*{Phase I reformulation}

Introduce the slack variable \(s \ge 0\) to account for infeasibility and modify the constraints to follow the convention used in the lecture:
\[
\begin{aligned}
\min_{x,y,s}\quad & s\\[0.3em]
\text{s.t.}\quad
& -x + 2y - 2 \le s,\\
& x + 2y - 6 \le 0,\\
& x - 2y - 2 \le 0,\\
& -x \le 0,\quad -y \le 0,\quad s \ge 0.
\end{aligned}
\]

Choose \(s_{0} = 2 = \max(0, g_1(x_{0}))\) so that Constraint (C1) becomes satisfied at the starting point:
\[
-1 + 2\cdot 2.5 - 2 = 2.
\]
Thus \((x_{0}, y_{0}, s_{0}) = (1,\, 2.5,\, 2)\) is feasible for the Phase I problem.

\subsection*{Solving Phase I}

Because the original problem is feasible, the minimum of Phase I is \(s^\star = 0\), attained at any point satisfying (C1)–(C5).

A suitable feasible point is found by intersecting (C1) and (C2):
\[
\begin{cases}
-x + 2y - 2 = 0,\\
x + 2y - 6 = 0.
\end{cases}
\]
Solve the system:
\[
x = 2y - 2,
\]
\[
(2y - 2) + 2y - 6 = 0
\;\Rightarrow\; 4y - 8 = 0
\;\Rightarrow\; y = 2,
\]
\[
x = 2\cdot 2 - 2 = 2.
\]

Thus the candidate point is:
\[
(x,y,s) = (2,2,0).
\]

Check feasibility:
\[
\begin{aligned}
\text{C1: } &2 - 2\cdot 2 + 2 = 0 \ge 0,\\
\text{C2: } &-2 - 2\cdot 2 + 6 = 0 \ge 0,\\
\text{C3: } &-2 + 2\cdot 2 + 2 = 4 \ge 0,\\
\text{C4: } &2 \ge 0,\qquad
\text{C5: } 2 \ge 0,\\
s &= 0.
\end{aligned}
\]

Thus, a valid feasible point for the original problem, produced by the Phase I method, is
\[
x_{\text{feasible}} = \begin{pmatrix}2\\[0.2em]2\end{pmatrix},
\]
which satisfies all original constraints and can be used to start the active-set method.

\section*{"Big M" method for initial feasible point (Bonus 0.5 pt)}

We introduce the variable
\[
z := \begin{pmatrix}x\\y\\s\end{pmatrix} \in \mathbb{R}^3
\]
and rewrite the Big M formulation as a quadratic program of the form
\[
\min_{z} \; \frac12 z^\top Q z + q^\top z + r
\quad\text{subject to}\quad
Bz \le b.
\]

The Big M objective is
\[
F_M(x,y,s) = (x-1)^2 + (y-2.5)^2 + 100\,s,
\]
so
\[
Q = \begin{pmatrix}
2 & 0 & 0\\
0 & 2 & 0\\
0 & 0 & 0
\end{pmatrix},\qquad
q = \begin{pmatrix}-2\\-5\\100\end{pmatrix},\qquad
r = 7.25.
\]

The constraints
\[
\begin{aligned}
&-x + 2y - 2 \le s,\\
&x + 2y - 6 \le 0,\\
&x - 2y - 2 \le 0,\\
&-x \le 0,\quad -y \le 0,\quad s \ge 0
\end{aligned}
\]
are written as
\[
Bz \le b,\quad B\in\mathbb{R}^{6\times 3},\; b\in\mathbb{R}^6,
\]
with rows
\[
B =
\begin{pmatrix}
-1 & 2  & -1\\
 1 & 2  &  0\\
 1 & -2 &  0\\
-1 & 0  &  0\\
 0 & -1 &  0\\
 0 & 0  & -1
\end{pmatrix},\qquad
b = \begin{pmatrix}2\\6\\2\\0\\0\\0\end{pmatrix}.
\]

The starting point is
\[
z_0 = \begin{pmatrix}1\\[0.1em]2.5\\[0.1em]2\end{pmatrix},
\]
for which
\[
Bz_0 - b = \begin{pmatrix}
0\\[0.1em]0\\[0.1em]-6\\[0.1em]-1\\[0.1em]-2.5\\[0.1em]-2
\end{pmatrix}\le 0,
\]
so \(z_0\) is feasible and constraints 1 and 2 are active:
\[
\mathcal{A}_0 = \{1,2\}.
\]

\subsection*{Active-set method for the Big M QP}

For a given active set \(\mathcal{A}_k\), let \(B_k\) and \(b_k\) collect the
corresponding rows of \(B\) and entries of \(b\). The active-set subproblem is
\[
\min_{z} \;\frac12 z^\top Q z + q^\top z + r
\quad\text{s.t.}\quad B_k z = b_k,
\]
whose KKT system is
\[
\begin{pmatrix}
Q & B_k^\top\\
B_k & 0
\end{pmatrix}
\begin{pmatrix}
z_{k+1}^{\text{cand}}\\[0.2em] \lambda_{k+1,\mathcal{A}_k}
\end{pmatrix}
=
\begin{pmatrix}
-q\\[0.2em] b_k
\end{pmatrix}.
\]
This gives a candidate point \(z_{k+1}^{\text{cand}}\); the search direction is
\(
p_k = z_{k+1}^{\text{cand}} - z_k
\),
and we choose a step length \(\alpha_k\in(0,1]\) so that
\(z_k + \alpha_k p_k\) remains feasible. If \(\alpha_k < 1\), a new constraint
becomes active. If \(p_k = 0\), we inspect the multipliers and possibly drop a
constraint with negative multiplier.

\paragraph{Iteration 0.}
Active set: \(\mathcal{A}_0 = \{1,2\}\), so
\[
B_0 = \begin{pmatrix}
-1 & 2 & -1\\
 1 & 2 &  0
\end{pmatrix},\qquad
b_0 = \begin{pmatrix}2\\6\end{pmatrix}.
\]
The KKT system
\[
\begin{pmatrix}
Q & B_0^\top\\
B_0 & 0
\end{pmatrix}
\begin{pmatrix}
z_1^{\text{cand}}\\ \lambda_{1,\mathcal{A}_0}
\end{pmatrix}
=
\begin{pmatrix}
-q\\ b_0
\end{pmatrix}
\]
yields
\[
z_1^{\text{cand}} = \begin{pmatrix}81\\[0.1em]-37.5\\[0.1em]-158\end{pmatrix},
\qquad
\lambda_{1,1} = 100,\;\lambda_{1,2} = -60.
\]
The search direction is
\[
p_0 = z_1^{\text{cand}} - z_0
    = \begin{pmatrix}80\\[0.1em]-40\\[0.1em]-160\end{pmatrix}.
\]

The inactive constraints along \(z(\alpha)=z_0+\alpha p_0\) satisfy
\[
\begin{pmatrix}
160\alpha-6\\[0.1em]-80\alpha-1\\[0.1em]40\alpha-2.5\\[0.1em]160\alpha-2
\end{pmatrix}\le 0,
\]
which gives
\[
\alpha_0 \le \min\!\left\{\frac{6}{160},\frac{2.5}{40},\frac{2}{160}\right\}
= \frac{1}{80}.
\]
The blocking constraint is 6, so \(\alpha_0 = 1/80\) and
\[
z_1 = z_0 + \alpha_0 p_0
    = \begin{pmatrix}2\\[0.1em]2\\[0.1em]0\end{pmatrix}.
\]
At \(z_1\),
\(
Bz_1 - b = (0,0,-4,-2,-2,0)^\top
\),
so
\[
\mathcal{A}_1 = \{1,2,6\}.
\]

\paragraph{Iteration 1.}
With \(\mathcal{A}_1 = \{1,2,6\}\),
\[
B_1 = \begin{pmatrix}
-1 & 2 & -1\\
 1 & 2 &  0\\
 0 & 0 & -1
\end{pmatrix},\qquad
b_1 = \begin{pmatrix}2\\6\\0\end{pmatrix}.
\]
Solving
\[
\begin{pmatrix}
Q & B_1^\top\\
B_1 & 0
\end{pmatrix}
\begin{pmatrix}
z_2^{\text{cand}}\\ \lambda_{2,\mathcal{A}_1}
\end{pmatrix}
=
\begin{pmatrix}
-q\\ b_1
\end{pmatrix}
\]
gives
\[
z_2^{\text{cand}} = \begin{pmatrix}2\\[0.1em]2\\[0.1em]0\end{pmatrix}
= z_1,\qquad
\lambda_{2,1} = \tfrac{5}{4},\;
\lambda_{2,2} = -\tfrac{3}{4},\;
\lambda_{2,6} = \tfrac{395}{4}.
\]
Hence \(p_1 = 0\). Since one multiplier is negative (\(\lambda_{2,2}<0\)),
constraint 2 is removed from the active set:
\[
\mathcal{A}_2 = \{1,6\},\qquad z_2 := z_1.
\]

\paragraph{Iteration 2.}
Now
\[
B_2 = \begin{pmatrix}
-1 & 2 & -1\\
 0 & 0 & -1
\end{pmatrix},\qquad
b_2 = \begin{pmatrix}2\\0\end{pmatrix}.
\]
The KKT system
\[
\begin{pmatrix}
Q & B_2^\top\\
B_2 & 0
\end{pmatrix}
\begin{pmatrix}
z_3^{\text{cand}}\\ \lambda_{3,\mathcal{A}_2}
\end{pmatrix}
=
\begin{pmatrix}
-q\\ b_2
\end{pmatrix}
\]
yields
\[
z_3^{\text{cand}} =
\begin{pmatrix}
\frac{7}{5}\\[0.1em]\frac{17}{10}\\[0.1em]0
\end{pmatrix},\qquad
\lambda_{3,1} = \frac{4}{5},\;
\lambda_{3,6} = \frac{496}{5}.
\]
The search direction is
\[
p_2 = z_3^{\text{cand}} - z_2
    = \begin{pmatrix}-\tfrac{3}{5}\\[0.1em]-\tfrac{3}{10}\\[0.1em]0\end{pmatrix}.
\]

Along \(z(\alpha)=z_2+\alpha p_2\), the inactive constraints satisfy
\[
\begin{pmatrix}
-\tfrac{6}{5}\alpha\\[0.1em]-4\\[0.1em]\tfrac{3}{5}\alpha-2\\[0.1em]\tfrac{3}{10}\alpha-2
\end{pmatrix}\le 0
\]
for all \(\alpha\in[0,1]\). Thus no new constraint blocks the step and we take
\(\alpha_2 = 1\):
\[
z_3 = z_2 + p_2 =
\begin{pmatrix}
\frac{7}{5}\\[0.1em]\frac{17}{10}\\[0.1em]0
\end{pmatrix}.
\]
At \(z_3\),
\(
Bz_3 - b = (0,-\tfrac{6}{5},-4,-\tfrac{7}{5},-\tfrac{17}{10},0)^\top
\),
so
\[
\mathcal{A}_3 = \{1,6\}.
\]

\paragraph{Iteration 3 (optimality).}
Solving the KKT system again with \(\mathcal{A}_3 = \{1,6\}\) returns
\[
z_4^{\text{cand}} = z_3,\qquad
\lambda_{4,1} = \frac{4}{5},\;
\lambda_{4,6} = \frac{496}{5},
\]
so \(p_3 = 0\) and all multipliers of active constraints are nonnegative.
Therefore the KKT conditions are satisfied and the active-set method terminates
at
\[
(x^\star,y^\star,s^\star) =
\left(\frac{7}{5},\,\frac{17}{10},\,0\right) = \left(1.4, 1.7, 0\right).
\]

This matches the solution we found in the first exercise as expected.

\end{document}
